% ============================================================
%  SECTION 10 (was 8): Why exchange rates matter for India   (Member 4)
%  Selvi box changed at merge to ₹90 -> ₹95 to match Member 1's ₹95.
% ============================================================
\section{Why exchange rates matter for India}\label{sec:matters}

You have already met the winners and losers from a change in the rupee \seesec{sec:appdep}. Here we follow those effects through the wider economy. Whether you gain or lose depends on whether you earn dollars, spend them or owe them.

Start with those who earn dollars. When the rupee weakens, every dollar an exporter earns turns into more rupees. The exporter can keep the dollar price unchanged and enjoy a bigger income, or cut the price abroad to win more orders. Either way, Indian goods become easier to sell, and we say that their \keyterm{export competitiveness}{The ability of a country's firms to sell goods and services abroad at prices that foreign buyers are willing to pay.} has improved. The gain is smaller than it looks for firms that import their own raw materials or machines, as the story of Selvi's unit shows.

\begin{casestudy}{A weaker rupee in Selvi's unit}
Selvi's unit in Tiruppur sells cotton T-shirts to a buyer in the United States, priced in dollars. Last year she shipped an order worth \$20,000 when the rupee was at ₹90 per dollar, and it brought in ₹18 lakh. This year the same order, at ₹95, brings in ₹19 lakh. The extra ₹1 lakh can pay for overtime wages, or let her offer a small discount to win the next order.

This is not just a classroom story. In early 2025, the Tiruppur Exporters' Association estimated that the rupee's fall during 2024 accounted for about 5 per cent of the growth in the town's garment exports (Business Standard, 2025a).

But Selvi's gain is not pure profit. Some of her dyes, accessories and machine parts are imported, and these now cost more. Her American buyer reads the same news, too, and may ask her to share the gain through a lower price. And if the rupee turns and strengthens, the gain disappears; that is exactly the risk her forward contract guards against \seesec{sec:forex}.
\figsource{The order values and exchange rates in this example are hypothetical.}
\end{casestudy}

Now consider those who spend dollars. India imports crude oil, gold, electronics, fertiliser and edible oil, almost all priced in dollars. A weaker rupee makes each of them dearer in rupees, even when the dollar price has not changed. Diesel runs the trucks that carry vegetables, grain and medicines, so costlier fuel spreads into the price of almost everything (Figure~\ref{fig:inflation}). A rise in prices that comes from dearer imports is called \keyterm{imported inflation}{A rise in the general price level caused by higher rupee prices of imported goods, often because the rupee has depreciated.}.

\begin{figure}[!htb]
  \centering
  \resizebox{\linewidth}{!}{% Figure (Section 10, Member 4): how a weaker rupee turns into imported inflation.
% Same look as Member 1's flow diagram (Figure 2).
\begin{tikzpicture}[
  font=\sffamily\small,
  chainbox/.style={fill=fxnavy!8, draw=fxnavy, line width=0.8pt, rounded corners=1.5mm,
               text width=24mm, minimum height=19mm, align=center, inner sep=3pt},
  last/.style={chainbox, fill=fxorange, draw=fxorange, text=white},
  flow/.style={-{Stealth[length=2.6mm,width=2.2mm]}, line width=1.3pt, fxnavy},
  node distance=5mm,
]
\node[chainbox] (a) {\textbf{Rupee weakens}\\[2pt]{\footnotesize ₹85 → ₹90 per US dollar}};
\node[chainbox, right=of a] (b) {\textbf{Imports cost more}\\[2pt]{\footnotesize crude oil, edible oil, electronics, fertiliser}};
\node[chainbox, right=of b] (c) {\textbf{Fuel prices rise}\\[2pt]{\footnotesize petrol, diesel and LPG}};
\node[chainbox, right=of c] (d) {\textbf{Costs rise}\\[2pt]{\footnotesize transport and production}};
\node[last, right=of d] (e) {\textbf{Prices rise}\\[2pt]{\footnotesize food and everyday goods: imported inflation}};
\foreach \u/\v in {a/b, b/c, c/d, d/e} \draw[flow] (\u) -- (\v);
\end{tikzpicture}
}
  \caption{How a weaker rupee becomes imported inflation.}
  \label{fig:inflation}
  {\raggedright\small\sffamily Read the chain from left to right. Each step takes time, so the full effect on prices can take several months to appear. The exchange rates shown are hypothetical.\par}
  \figsource{Drawn by the authors.}
\end{figure}

\begin{workedexample}{The same barrel, a different price}
Suppose a barrel of crude oil costs \$80 on the world market (hypothetical figures).
\begin{itemize}
  \item At ₹85 per US dollar, India pays 80 × 85 = ₹6,800 per barrel.
  \item At ₹90 per US dollar, India pays 80 × 90 = ₹7,200 per barrel.
\end{itemize}
The oil has not become any dearer in dollars, yet India pays ₹400 more for every barrel, a rise of about 5.9 per cent.

\textbf{Now you try.} The dollar price of oil rises to \$90 while the rupee stays at ₹90 per dollar. What is the new rupee price of a barrel, and by what percentage has it risen compared with ₹6,800?
\answers{₹8,100, a rise of about 19 per cent. When oil prices rise and the rupee weakens together, the two effects multiply.}
\end{workedexample}

Then there are those who owe dollars. A company that borrowed in dollars, such as an airline paying for new aircraft, must now find more rupees to repay the same loan. A sharp fall in the rupee can turn a healthy balance sheet into a worrying one.

All of this matters to the Reserve Bank of India. When the RBI raises its policy interest rate (the repo rate), Indian assets pay more, foreign money tends to flow in and the rupee strengthens. Imports become cheaper and inflation cools, although exporters find life harder. This route from interest rates to spending and prices, working through the exchange rate, is called the \keyterm{exchange rate channel}{The way a change in the central bank's policy interest rate affects spending and prices by moving the exchange rate, and through it exports and imports.} of monetary policy (Jayadev, 2021, Section 15.9). It is one reason the RBI keeps a close eye on the rupee whenever it sets interest rates.
